% ============================================================
\chapter*{Pr\'eface}
\addcontentsline{toc}{chapter}{Pr\'eface}
\markboth{Pr\'eface}{Pr\'eface}
% ============================================================

\section*{Objectifs de ce cours}

La statistique bay\'esienne constitue l'un des deux grands paradigmes
de l'inf\'erence statistique, aux c\^ot\'es de l'approche fr\'equentiste.
Fond\'ee sur le th\'eor\`eme de Bayes, elle offre un cadre coh\'erent
pour quantifier l'incertitude, incorporer des connaissances a priori
et mettre \`a jour les croyances \`a la lumi\`ere des donn\'ees.

Ce cours s'adresse aux \'etudiants de niveau Master en math\'ematiques
appliqu\'ees, en science des donn\'ees et en intelligence artificielle.
Il suppose une ma\^itrise pr\'ealable de la th\'eorie des probabilit\'es,
de la statistique math\'ematique \'el\'ementaire et de l'alg\`ebre lin\'eaire.

\section*{Organisation du cours}

Le cours est structur\'e en onze chapitres, regroup\'es en quatre parties
th\'ematiques~:

\begin{enumerate}
  \item \textbf{Fondements} (Chapitres 1--5)~: paradigme bay\'esien,
        familles conjugu\'ees, estimation, tests et pr\'ediction.
  \item \textbf{Mod\`eles avanc\'es} (Chapitre 6)~: mod\`eles hi\'erarchiques.
  \item \textbf{M\'ethodes de calcul} (Chapitres 7--9)~: MCMC
        (Metropolis--Hastings, \'echantillonneur de Gibbs) et inf\'erence
        variationnelle.
  \item \textbf{Extensions et applications} (Chapitres 10--11)~:
        statistique bay\'esienne non param\'etrique et applications
        concr\`etes.
\end{enumerate}

\section*{Philosophie p\'edagogique}

Chaque chapitre alterne entre~:
\begin{itemize}
  \item des \textbf{d\'eveloppements th\'eoriques} rigoureux (d\'efinitions,
        th\'eor\`emes, preuves),
  \item des \textbf{exemples d\'etaill\'es} avec calculs explicites,
  \item des \textbf{impl\'ementations num\'eriques} en Python (PyMC, NumPy,
        SciPy),
  \item des \textbf{exercices} de difficult\'e croissante.
\end{itemize}

\section*{Approche bay\'esienne \textit{versus} fr\'equentiste}

Avant d'entamer l'\'etude d\'etaill\'ee, situons les deux paradigmes~:

\begin{center}
\renewcommand{\arraystretch}{1.3}
\begin{tabular}{p{4cm} p{5cm} p{5cm}}
\toprule
\textbf{Aspect} & \textbf{Fr\'equentiste} & \textbf{Bay\'esien} \\
\midrule
Param\`etre $\theta$ & Fixe, inconnu & Variable al\'eatoire \\
Probabilit\'e & Fr\'equence limite & Degr\'e de croyance \\
Incertitude & Intervalles de confiance & Intervalles de cr\'edibilit\'e \\
Information a priori & Non utilis\'ee formellement & Int\'egr\'ee via le prior \\
R\'esultat principal & Estimateur $\hat{\theta}(x)$ & Loi a posteriori $\pi(\theta \mid x)$ \\
Coh\'erence & Peut violer le principe de vraisemblance & Coh\'erence assur\'ee \\
Petit \'echantillon & Difficult\'es possibles & Naturellement adapt\'ee \\
Calcul & Souvent analytique & Souvent num\'erique (MCMC) \\
\bottomrule
\end{tabular}
\end{center}

\section*{Le th\'eor\`eme fondateur}

Tout le paradigme bay\'esien repose sur la formule~:

\begin{formules}[title=\textbf{Th\'eor\`eme de Bayes}]
Soit $\theta \in \Theta$ un param\`etre de loi a priori $\pi(\theta)$
et $x = (x_1, \ldots, x_n)$ un \'echantillon de vraisemblance
$L(\theta \mid x)$. La \textbf{loi a posteriori} est~:
\[
  \pi(\theta \mid x)
  = \frac{L(\theta \mid x)\,\pi(\theta)}
         {\int_{\Theta} L(\theta \mid x)\,\pi(\theta)\,d\theta}
  \propto L(\theta \mid x)\,\pi(\theta).
\]
\end{formules}

\section*{Notation et conventions}

Tout au long de ce cours, nous adoptons les conventions suivantes~:

\begin{itemize}
  \item $X$ d\'esigne une variable al\'eatoire, $x$ une r\'ealisation.
  \item $\theta$ est le param\`etre d'int\'er\^et, $\Theta$ l'espace
        des param\`etres.
  \item $\pi(\theta)$ est la loi a priori (\textit{prior}).
  \item $\pi(\theta \mid x)$ est la loi a posteriori (\textit{posterior}).
  \item $L(\theta \mid x) = f(x \mid \theta)$ est la vraisemblance.
  \item $m(x) = \int L(\theta \mid x)\,\pi(\theta)\,d\theta$ est la
        vraisemblance marginale (\textit{evidence}).
  \item $\E[\cdot]$, $\Var[\cdot]$, $\Cov[\cdot]$ d\'esignent
        esp\'erance, variance et covariance.
  \item $\PP(\cdot)$ d\'esigne la probabilit\'e.
  \item $\KL(p \| q)$ d\'esigne la divergence de Kullback--Leibler.
  \item $\norm{\cdot}$ d\'esigne la norme euclidienne.
\end{itemize}

\section*{Distributions utilis\'ees fr\'equemment}

\begin{center}
\renewcommand{\arraystretch}{1.3}
\begin{tabular}{lll}
\toprule
\textbf{Distribution} & \textbf{Notation} & \textbf{Param\`etres} \\
\midrule
Bernoulli & $\text{Ber}(p)$ & $p \in [0,1]$ \\
Binomiale & $\text{Bin}(n, p)$ & $n \in \N$, $p \in [0,1]$ \\
Poisson & $\text{Poi}(\lambda)$ & $\lambda > 0$ \\
Normale & $\mathcal{N}(\mu, \sigma^2)$ & $\mu \in \R$, $\sigma^2 > 0$ \\
Gamma & $\text{Ga}(\alpha, \beta)$ & $\alpha, \beta > 0$ \\
B\^eta & $\text{Be}(a, b)$ & $a, b > 0$ \\
Dirichlet & $\text{Dir}(\boldsymbol{\alpha})$ & $\alpha_k > 0$ \\
Student & $t_\nu(\mu, \sigma^2)$ & $\nu > 0$ \\
Wishart & $\mathcal{W}_p(\nu, \boldsymbol{\Sigma})$ & $\nu > p-1$ \\
Inverse-Gamma & $\text{IG}(\alpha, \beta)$ & $\alpha, \beta > 0$ \\
\bottomrule
\end{tabular}
\end{center}

\section*{Jalons historiques}

\begin{itemize}
  \item \textbf{1763}~: Publication posthume du th\'eor\`eme de Thomas Bayes
        par Richard Price.
  \item \textbf{1812}~: Pierre-Simon de Laplace formalise et g\'en\'eralise
        la m\'ethode dans \textit{Th\'eorie analytique des probabilit\'es}.
  \item \textbf{1937}~: Bruno de Finetti d\'eveloppe la notion de
        probabilit\'e subjective et le th\'eor\`eme d'\'echangeabilit\'e.
  \item \textbf{1954}~: Leonard Jimmie Savage publie \textit{The Foundations
        of Statistics}, pilier de la th\'eorie de la d\'ecision bay\'esienne.
  \item \textbf{1970s}~: Dennis Lindley et d'autres \'etablissent les
        fondements modernes.
  \item \textbf{1990s}~: R\'evolution MCMC --- Metropolis--Hastings devient
        le cheval de bataille du calcul bay\'esien.
  \item \textbf{2000s--2020s}~: Inf\'erence variationnelle, processus
        gaussiens, apprentissage profond bay\'esien.
\end{itemize}

\section*{Logiciels et biblioth\`eques}

Les impl\'ementations pratiques de ce cours utilisent~:

\begin{itemize}
  \item \textbf{Python 3.10+} comme langage principal.
  \item \textbf{PyMC} (v5+) pour la mod\'elisation bay\'esienne et MCMC.
  \item \textbf{ArviZ} pour le diagnostic et la visualisation des cha\^ines.
  \item \textbf{NumPy / SciPy} pour le calcul num\'erique.
  \item \textbf{Matplotlib / Seaborn} pour les graphiques.
\end{itemize}

\begin{codeblock}[title=\textbf{Installation des d\'ependances}]
\begin{minted}{python}
# Installation recommandée
# pip install pymc arviz numpy scipy matplotlib seaborn

import pymc as pm
import arviz as az
import numpy as np
import matplotlib.pyplot as plt

print(f"PyMC version: {pm.__version__}")
print(f"ArviZ version: {az.__version__}")
\end{minted}
\end{codeblock}

\section*{Comment utiliser ce cours}

\begin{attention}[title=\textbf{Pr\'erequis importants}]
Ce cours suppose que l'\'etudiant ma\^itrise~:
\begin{enumerate}
  \item La th\'eorie des probabilit\'es (variables al\'eatoires, lois
        classiques, convergences).
  \item La statistique math\'ematique (estimation, tests, vraisemblance).
  \item L'alg\`ebre lin\'eaire (matrices, d\'ecompositions, formes
        quadratiques).
  \item La programmation en Python (syntaxe de base, NumPy).
\end{enumerate}
\end{attention}

\begin{intuition}[title=\textbf{Fil conducteur}]
L'id\'ee directrice de ce cours est simple~: \textit{toute inf\'erence
est un probl\`eme de mise \`a jour d'une distribution de probabilit\'e}.
On part d'une croyance a priori $\pi(\theta)$, on observe des donn\'ees
$x$, et on obtient une croyance a posteriori $\pi(\theta \mid x)$.
Cette vision unifi\'ee sous-tend l'ensemble des m\'ethodes pr\'esent\'ees.
\end{intuition}

\section*{R\'ef\'erences principales}

\begin{itemize}
  \item A.~Gelman, J.B.~Carlin, H.S.~Stern, D.B.~Dunson, A.~Vehtari,
        D.B.~Rubin, \textit{Bayesian Data Analysis} (3e \'ed.), 2013.
  \item C.P.~Robert, \textit{The Bayesian Choice} (2e \'ed.), 2007.
  \item J.K.~Kruschke, \textit{Doing Bayesian Data Analysis} (2e \'ed.),
        2015.
  \item D.~Barber, \textit{Bayesian Reasoning and Machine Learning}, 2012.
  \item K.P.~Murphy, \textit{Machine Learning: A Probabilistic Perspective},
        2012.
  \item P.~Hoff, \textit{A First Course in Bayesian Statistical Methods},
        2009.
\end{itemize}

\bigskip
\begin{flushright}
\textit{L'auteur} \\
Mars 2026
\end{flushright}
